\section{Материальное обеспечение группы}
\begin{table}[h!]
	\centering
		\renewcommand{\arraystretch}{1.25}
			\setlength{\tabcolsep}{10pt}
				\begin{tabular}{|p{7.2cm}|c|p{7.2cm}|c|}
					\hline
					\multicolumn{4}{|c|}{\textbf{Специальное снаряжение}} \\
					\hline
					\multicolumn{2}{|c|}{\textbf{Групповое}} & \multicolumn{2}{c|}{\textbf{Личное}} \\
					\hline
					\textbf{Наименование} & \textbf{Кол-во} & \textbf{Наименование} & \textbf{Кол-во} \\
					\hline
					GPS + комплект карт & 2 & Карта-схема маршрута & 1 \\
					\hline
					Палатка «Зима» & 1 & Компас & 1 \\
					\hline
					Печь варочная с канами & 1 & Костюм ветрозащитный & 1 \\
					\hline
					Пила & 1 & Ветрозащитная маска & 1 \\
					\hline
					Топор & 1 & Лыжи с тросиковыми креплениями + палки & 1 \\
					\hline
					Спас конец & 1 & Бахилы & 1 \\
					\hline
					Лопата & 2 & Лавинный щуп & 1 \\
					\hline
				\end{tabular}
	\end{table}
\begin{table}[h!]

Палатка была взята лёгкая (4.5 кг вместе с налипшим льдом) - экономия веса засчёт оптимизации фурнитуры и отсутствии тента. Вся ткань была пропитана силиконом, однако в этом было больше минусов, чем плюсов: несмотря на полную защиту от ветра, еженочно на участников капало увеличенное количество конденсата, сушиться приходилось серьёзнее. 
Печка тоже лёгкая (2.8 кг), варочная, производства одного из наших участников (Михаила Чернецкого), с увеличенной тягой и отдельными изолирующими ножками (поэтому не пришлось дополнительно таскать лавлист). Каны брали три штуки (чай, еда и термоса), тоже облегчённые и тоже производства Михаила. 
На обеды было взято 6 термосов фирмы Арктика (по 1.2 литра каждый), половина для супов, половина для чая.
Крепления на лыжах участников были не слишлом разнообразными: "Азимут" и с качающимися щёчками. Все остались целыми к концу похода.
Рвалось несколько тросиков, но у каждого участника было 1-2 запасномых тросика, поэтому проблем не было.
Все участники шли на камусах, кроме Михаила Чернецкого (он был на лыжах с насечками). 5 участников использовали камуса фирмы Маяк, 3 человека - фирмы Contour. Проблема маяковских камусов была в основном в креплении на носке лыжи: камус со слишком большой скобой был соединена тугой резинкой, а сами скобы были несварены и не замкнуты в целом, они слетали. Камуса фирмы Contour показали себя хорошо (но один из них случайно пожгли углями с печки).
При восхождении к Ицылу был потерян бегунок на "лягушке" (так как лыжи были в рюкзаке, болтались крепления, а стопор для бегунка отсутствовал, поэтому последний раскрутился и затерялся среди камней), но кустарным методом уже после спуска была придумана альтернатива из болта и дугообразной железки (пожертвовали постояльцы в избе, в которой остановились на обед). 
Санки не брали, так как мы ожидали большую часть лесной зоны и курумников, где использование саней неудобно, да и для первой категории маршрута нецелесообразно.
Газа взяли 5 баллонов по 220 грамм, на случай холодной ночёвки или необходимости экономить дрова. Пригодились почти все баллоны (так как не сильно везло с дровами).
Средний вес общественного снаряжения для парней был 3.8 кг, у девушек 3.2 кг, шлось хорошо.
Средний вес общественного снаряжения - 15 кг на человека.
Продукты и топливо (газ в цанговых баллонах) - 0.722 кг в день/4.32 кг на все дни на человека.

\clearpage
